% C015 椭圆焦点三角形四心坐标性质——典型例题与变式
\begin{examplebox}
\textbf{例 1（四心同求：选一个可以心算的点）}

椭圆 $\dfrac{x^2}{25}+\dfrac{y^2}{16}=1$，
$P=(3,\tfrac{16}{5})$。求 $\triangle PF_1F_2$ 的四心，
并解释垂心为什么落在焦点处。

\textbf{解：} $a=5,b=4,c=3,e=3/5$，于是
\[
 G=\left(1,\frac{16}{15}\right),\qquad
 I=\left(\frac95,\frac65\right),
\]
\[
 H=(3,0),\qquad O'=\left(0,\frac85\right).
\]
检验内心：$PF_1=5+\frac35\cdot3=\frac{34}{5}$，
$PF_2=5-\frac35\cdot3=\frac{16}{5}$。
因此
\[
 x_I=\frac{PF_1-PF_2}{2}=\frac95,\qquad
 r=\frac{c|y_0|}{a+c}=\frac65.
\]
\textbf{几何解释：} $P$ 的横坐标恰为 $c=3$，
所以 $PF_2$ 竖直，$F_1F_2$ 水平，$\angle PF_2F_1=90^\circ$。
直角三角形的垂心就是直角顶点 $F_2=(3,0)$。
同时 $P-2O'=(3,16/5)-(0,16/5)=(3,0)$，交叉验证成立。

\tcbline
\textbf{例 2（秒求重心、内心轨迹）}

同一椭圆 $\dfrac{x^2}{25}+\dfrac{y^2}{16}=1$，
$P$ 在椭圆上运动，但不能取 $(\pm5,0)$。
分别求重心 $G=(X,Y)$ 和内心 $I=(X,Y)$ 的轨迹。

\textbf{解：}

重心 $X=x_0/3$、$Y=y_0/3$，反解为 $x_0=3X$、$y_0=3Y$。
代回椭圆：
\[
 \boxed{\frac{X^2}{25/9}+\frac{Y^2}{16/9}=1},
 \qquad (X,Y)\ne\left(\pm\frac53,0\right).
\]

内心 $X=\frac35x_0$、$Y=\frac38y_0$，
反解为 $x_0=\frac53X$、$y_0=\frac83Y$。
代回椭圆：
\[
 \boxed{\frac{X^2}{9}+\frac{Y^2}{9/4}=1},
 \qquad (X,Y)\ne(\pm3,0).
\]
\textbf{认识迁移：} 两条轨迹都用\textbf{设新点 $\to$ 反解旧点 $\to$ 代回原方程}，
无须从头求四心。

\tcbline
\textbf{例 3（短轴端点：四心共线，却未必重合）}

在椭圆 $\dfrac{x^2}{25}+\dfrac{y^2}{16}=1$ 上取 $P=(0,4)$。

\textbf{解：} 三角形关于 $y$ 轴对称，故四心均在 $y$ 轴上：
\[
 G=\left(0,\frac43\right),\qquad
 I=\left(0,\frac32\right),
\]
\[
 H=\left(0,\frac94\right),\qquad
 O'=\left(0,\frac78\right).
\]
\textbf{关键辨析：} 等腰三角形的四心只保证共轴，
\textbf{不保证重合}。它们是否重合还要看是不是等边三角形。

\tcbline
\textbf{例 4（外心竟在底边的另一边？）}

椭圆 $\dfrac{x^2}{25}+\dfrac{y^2}{9}=1$，
焦点为 $(\pm4,0)$，取 $P=(0,3)$。求外心并解释其符号。

\textbf{解：}
\[
 O'=\left(0,\frac{0+9-16}{2\cdot3}\right)
 =\boxed{\left(0,-\frac76\right)}.
\]
$P$ 位于 $x$ 轴上方，但外心在下方；这是完全可能的。
因为 $OP=3<c=4$，点 $P$ 在以 $F_1F_2$ 为直径的圆内，
故 $\angle F_1PF_2>90^\circ$，三角形顶角为钝角，
外心在顶点 $P$ 所在半平面的另一侧。
本题垂心为 $H=(0,16/3)$，而 $P-2O'=(0,3+7/3)=(0,16/3)$。

\tcbline
\textbf{例 5（逆向应用：已知内心还原椭圆动点）}

椭圆 $\dfrac{x^2}{25}+\dfrac{y^2}{16}=1$ 的焦点三角形
内心为 $I=(\tfrac95,\tfrac65)$，求 $P$ 与三角形面积。

\textbf{解：} 先反解坐标，
\[
 x_0=\frac{x_I}{e}=\frac{9/5}{3/5}=3,
 \qquad
 y_0=\frac{a+c}{c}y_I=\frac83\cdot\frac65=\frac{16}{5}.
\]
故 $P=(3,16/5)$。检验：$9/25+(16/5)^2/16=9/25+16/25=1$。
面积为
\[
 \boxed{S=c|y_0|=3\cdot\frac{16}{5}=\frac{48}{5}}.
\]
\textbf{提示：} 逆推得到的 $P$ 必须回代椭圆方程；
并非任意平面点都可能是该椭圆某个焦点三角形的内心。

\tcbline
\textbf{例 6（高价值变式：什么时候外心是椭圆中心？）}

设 $P$ 为椭圆上的非长轴端点，判断何时 $O'=O=(0,0)$，
并找出椭圆离心率的条件。

\textbf{解：} 外心在原点当且仅当
\[
 x_0^2+y_0^2=c^2.
\]
即 $P$ 同时落在椭圆和以 $F_1F_2$ 为直径、半径为 $c$ 的圆上。
因为椭圆上 $OP^2$ 的最小值是 $b^2$，最大值是 $a^2$，
而 $c<a$，两曲线存在交点当且仅当 $c\ge b$。
化成离心率：
\[
 c^2\ge b^2=a^2-c^2
 \ \Longleftrightarrow\ 2c^2\ge a^2
 \ \Longleftrightarrow\
 \boxed{e\ge\frac{1}{\sqrt2}}.
\]
当 $e=1/\sqrt2$ 时，只有 $P=(0,\pm b)$ 满足；
当 $e>1/\sqrt2$ 时，有上下半平面中关于坐标轴对称的四个交点。
这些位置的焦点三角形顶角均为直角。

\textbf{变式提醒：} 题目出现“外心恰为椭圆中心”，
立刻联想\textbf{直径所对的圆周角}，比强行求外心更直观。
\end{examplebox}
